%! AP Statistics — Unit 3, Lesson 3.2 — Student Guided Notes
\documentclass[10pt]{article}
\usepackage{apstats-article}
\usepackage{apstats-boxes}

\begin{document}

\pageheader{Unit 3, Lesson 3.2}{Guided Notes}
\namedateperiod

\vspace{0.12in}

\begin{objectivebox}
\small By the end of this lesson, I will be able to\dots\\[4pt]
\writeline\\[1pt]
\writeline\\[1pt]
\writeline
\end{objectivebox}

\vspace{0.1in}

\begin{vocabbox}
\termblanklong{Population}
\termblanklong{Sample}
\termblanklong{Observational study}
\termblanklong{Experiment}
\termblanklong{Generalizability}
\termblanklong{Causal relationship}
\end{vocabbox}

\vspace{0.1in}

\begin{hookbox}
\small\textbf{Two headlines:}
\begin{itemize}[itemsep=2pt]
  \item \textit{``Study: People who own dogs live longer.''}
  \item \textit{``Clinical Trial: New drug extends life expectancy by 3 years.''}
\end{itemize}
\textbf{Which headline tells us the variable \emph{caused} the effect? How do you know?}\\[4pt]
\writeline\\[1pt]\writeline
\end{hookbox}

\vspace{0.1in}

\begin{notesbox}{Population vs. Sample}
\small
\renewcommand{\arraystretch}{1.6}
\begin{tabularx}{\linewidth}{@{} >{\bfseries\color{navy}}p{0.22\linewidth} X @{}}
Population & All \blank{5cm} of interest. \\
Sample & A \blank{3cm} of the population selected for study. \\
\end{tabularx}

\vspace{10pt}
\textbf{Example:} A researcher surveys 500 randomly chosen U.S. adults about their
exercise habits.
\begin{itemize}[itemsep=2pt]
  \item Population: \blank{9cm}
  \item Sample: \blank{9cm}
\end{itemize}
\end{notesbox}

\vspace{0.1in}

\begin{notesbox}{Observational Study vs. Experiment}
\small
\renewcommand{\arraystretch}{1.6}
\begin{tabularx}{\linewidth}{@{} >{\bfseries\color{navy}}p{0.22\linewidth} X @{}}
Observational study & Researchers \blank{3cm} without \blank{4cm} treatments.
Includes surveys, retrospective and prospective studies. \\[4pt]
Experiment & Researchers \blank{3cm} treatments to \blank{4cm} units. \\
\end{tabularx}

\vspace{10pt}
\textbf{Conclusion Rules (DAT-2.B):}
\begin{enumerate}[itemsep=4pt]
  \item Can only generalize to the \blank{5cm} from which the sample was \blank{3cm} selected.
  \item \textbf{Observational study} $\to$ \blank{4cm} establish causation (DAT-2.B.3).
  \item \textbf{Experiment with random assignment} $\to$ \blank{4cm} establish causation.
\end{enumerate}
\end{notesbox}

\vspace{0.1in}

\begin{practicebox}
\small Classify each study. Then answer: can we generalize? Can we conclude causation?

\vspace{6pt}
\renewcommand{\arraystretch}{1.8}
\begin{tabularx}{\linewidth}{@{} X p{2.4cm} p{2cm} p{2cm} @{}}
\textbf{Study} & \textbf{Type} & \textbf{Generalize?} & \textbf{Causation?} \\
\hline
Researchers randomly assign 100 patients to drug or placebo; measure recovery time. &
\blank{2.2cm} & Yes\ \ \ No & Yes\ \ \ No \\
A survey of 500 randomly selected adults asks about screen time and sleep quality. &
\blank{2.2cm} & Yes\ \ \ No & Yes\ \ \ No \\
A teacher asks the 30 students in her class whether they like the new homework policy. &
\blank{2.2cm} & Yes\ \ \ No & Yes\ \ \ No \\
\end{tabularx}
\end{practicebox}

\end{document}
